\chapter{Excertos de Código e Configurações}\label{ap:codigo}
\thispagestyle{fancy}

Neste anexo, apresentam-se alguns excertos de código e configurações críticas desenvolvidas para o ecossistema tecnológico, evidenciando as decisões de integração e segurança tomadas durante o projeto.

\section{Configuração do Proxy Reversível em Next.js}\label{ap:proxy}

De forma a mascarar a localização real da API central e permitir a transmissão segura de cookies com a diretiva \textit{HttpOnly}, configurou-se uma regra de reencaminhamento no ficheiro de configuração do Next.js. O trecho abaixo ilustra como os pedidos direcionados para \texttt{/api/v1/*} são mapeados internamente para a API em .NET:

\begin{lstlisting}[language=HTML, caption={Configuração de rewrites de rede em Next.js.}]
// next.config.mjs
const nextConfig = {
  async rewrites() {
    return [
      {
        source: '/api/v1/:path*',
        destination: process.env.NEXT_PUBLIC_API_URL + '/api/v1/:path*',
      },
    ];
  },
};

export default nextConfig;
\end{lstlisting}

\section{Configuração de Políticas CORS no .NET API}\label{ap:cors}

Para garantir que apenas os clientes autorizados conseguem obter dados da API, implementou-se uma política de CORS estrita na inicialização do serviço .NET, limitando o acesso ao domínio do portal cliente:

\begin{lstlisting}[language=HTML, caption={Configuração de políticas CORS na API em C\#.}]
// Program.cs
var builder = WebApplication.CreateBuilder(args);

builder.Services.AddCors(options =>
{
    options.AddPolicy("PortalClientePolicy", policy =>
    {
        policy.WithOrigins(builder.Configuration["AllowedOrigins:PortalCliente"] 
                ?? "https://portal.20recolher.pt")
              .AllowAnyHeader()
              .AllowAnyMethod()
              .AllowCredentials(); // Permitir envio de tokens de cookies
    });
});

var app = builder.Build();
app.UseCors("PortalClientePolicy");
\end{lstlisting}

\section{Mecanismo de Filtro Global de Soft-Delete no EF Core}\label{ap:softdelete}

Para automatizar a exclusão lógica ao nível da base de dados sem necessidade de adicionar cláusulas condicionais repetitivas nas consultas do Entity Framework Core, configurou-se um filtro de consulta global no contexto de dados:

\begin{lstlisting}[language=HTML, caption={Configuração de filtro global de consulta para soft-delete no EF Core.}]
// ApplicationDbContext.cs
protected override void OnModelCreating(ModelBuilder modelBuilder)
{
    base.OnModelCreating(modelBuilder);

    // Aplicar filtro de eliminação lógica para todas as entidades que herdam de ISoftDeletable
    foreach (var entityType in modelBuilder.Model.GetEntityTypes())
    {
        if (typeof(ISoftDeletable).IsAssignableFrom(entityType.ClrType))
        {
            modelBuilder.Entity(entityType.ClrType)
                .HasQueryFilter(ConvertFilterExpression(entityType.ClrType));
        }
    }
}
\end{lstlisting}

\chapter{Registo Fotográfico do Lab de Hardware}\label{ap:fotos}
\thispagestyle{fancy}

Neste anexo, apresentam-se fotografias do espaço de trabalho e dos processos manuais de diagnóstico e manutenção efetuados no departamento de hardware da empresa.

\begin{figure}[H]
    \centering
    \confidencial % Indica que as fotografias são de cariz operacional privado
    \vspace{2cm}
    \textbf{[FOTOGRAFIA 1: Área de Triage e Diagnóstico de REEE]} \\
    \textit{Espaço reservado aos testes de RAM, discos e transformadores.}
    \vspace{2cm}
    \caption{Bancada de ensaio e computadores de teste utilizados no laboratório.}
    \label{fig:foto_bancada}
\end{figure}

\begin{figure}[H]
    \centering
    \confidencial
    \vspace{2cm}
    \textbf{[FOTOGRAFIA 2: Processo de Recondicionamento de Equipamentos]} \\
    \textit{Etapa de higienização interna e aplicação de composto térmico na CPU.}
    \vspace{2cm}
    \caption{Manutenção física e preparação estética de portátil recondicionado.}
    \label{fig:foto_manutencao}
\end{figure}
